\chapter{INTRODUCTION}

The development of the modular ZK-rollup benchmarking framework required coordinated implementation across multiple system layers, including workload generation, transaction sequencing, execution, proof handling, Layer-1 settlement, data availability submission, and post-experiment metric analysis. While the main final report presents the collective architecture, experimental methodology, and research findings of the project, this Individual Contribution Report focuses specifically on the implementation work, integration responsibilities, and code-level changes that I contributed to the project.

My contribution was concentrated on the outer and evaluation-facing layers of the rollup pipeline. In particular, I worked on the components that connect finalized Layer-2 batches to the Layer-1 settlement environment and on the tooling required to execute repeatable benchmark experiments. This included the Submitter service, which is responsible for preparing finalized batch payloads, selecting the configured data availability strategy, submitting commitments and metadata to the Layer-1 bridge contract, waiting for transaction confirmation, and exporting submission-related metrics. I also contributed to the supporting smart contract layer, including bridge and data availability related contract components used by the local benchmark environment.

In addition to the settlement path, I contributed to the benchmarking infrastructure that made the prototype suitable for empirical evaluation. This work included the initial benchmark suite, experiment configuration files, workload generation scripts, benchmark orchestration scripts, service readiness handling, and Python-based data tools for aggregating and analysing raw benchmark outputs. These components were necessary to transform the prototype from a set of independent services into a measurable end-to-end research pipeline. The implemented flow connects the workload generator, sequencer, executor, prover, submitter, Layer-1 contracts, and metric-processing tools into a reproducible benchmark workflow.

A significant part of my work also involved debugging and consistency fixes across the integration boundary between services. These fixes included improving Submitter reliability during Layer-1 confirmation handling, correcting configuration and parsing issues that affected benchmark execution, and ensuring that metrics produced by different services could be consistently consumed by the data analysis pipeline. Such fixes were important because incorrect settlement status handling, unstable service startup, or inconsistent metric schemas could directly affect the validity of the experimental results.

Since the project was developed collaboratively, this report does not simply describe the overall system. Instead, it traces my individual contribution using concrete software evidence. The analysis is based on Git commit history, commit diffs, file-level changes, source code inspection, and the final implementation structure. This approach allows the report to distinguish between shared project outcomes and the specific modules, refactors, bug fixes, and integration work that can be attributed to my contribution. Therefore, the following chapters present both a technical explanation of the code I implemented and an evidence-based mapping between my commits and the resulting project functionality.

